\section{Wprowadzenie}\label{sec:wprowadzenie}

Ten dokument opisuje framework letsgolegacy w całości: problem, który rozwiązuje,
architekturę, każdy komponent i każde repozytorium, formaty plików, polecenia oraz
procedurę pracy krok po kroku. Jest przeznaczony dla inżyniera, który ma z frameworkiem
pracować, rozwijać go albo ocenić jego wynik. Opisuje stan repozytoriów na konkretnych
commitach (sekcja~\ref{sec:wprowadzenie-wersja}) i niczego do nich nie dodaje.

\subsection{Problem}\label{sec:wprowadzenie-problem}

Systemy biznesowe na .NET Framework są przenoszone na współczesny .NET coraz częściej
przez agentów migracyjnych, takich jak modernize-dotnet (GitHub Copilot app
modernization for .NET). Kryterium odbioru, którym posługuje się samo narzędzie, to:
kandydat się buduje, a istniejące testy przechodzą. To kryterium nie mówi, czy system po
migracji zachowuje się tak samo wobec użytkownika. Migracja z .NET Framework na .NET
zmienia na przykład bibliotekę globalizacji z NLS na ICU, więc porównywanie i sortowanie
napisów zależne od kultury może dać inny wynik, choć build i testy są zielone. Tę
pułapkę bench nopCommerce traktuje jako główny kandydat na moment, wokół którego zbudowana
jest demonstracja: agent migruje, build i testy przechodzą, a weryfikacja i tak znajduje
realną różnicę zachowania.

Rdzeń opisuje swoje zadanie jednym zdaniem: niezależny, deterministyczny dowód, że system
.NET zmigrowany przez agenta AI nadal dotrzymuje kontraktu biznesowego. Nazwa Second Key
pochodzi z obrazu dwóch kluczy. Agent migracyjny przekręca pierwszy klucz: migruje kod.
Second Key przekręca drugi: nagrywa, jak obecny system naprawdę się zachowuje, odtwarza to
zachowanie na wersji zmigrowanej, sprawdza każdą różnicę względem spisanego kontraktu
tego, co musi się dziać i czego nie wolno, i daje osobie zatwierdzającej evidence pack
(pakiet dowodowy) dla każdej zmiany. \finding{Second Key weryfikuje; niczego nie
przepisuje.}

\zrodla{\path{letsgolegacy.secondkey/README.md},
\path{letsgolegacy.secondkey-nopcommerce-bench/docs/WORKPLAN.md},
\path{letsgolegacy.secondkey-nopcommerce-bench/docs/P6-RUNBOOK.md}}

\subsection{Czym jest framework letsgolegacy}\label{sec:wprowadzenie-framework}

Framework letsgolegacy to metoda weryfikowalnej modernizacji systemów legacy .NET razem z
narzędziami, które ją realizują. Narzędzia mieszkają w repozytoriach o prefiksie
\repopath{letsgolegacy.}:

\begin{itemize}
  \item \textbf{Second Key} --- łańcuch weryfikacji: nagranie ruchu (capture), kontrakt,
        odtworzenie (replay) na obu wersjach, porównanie (compare), werdykt i pakiet
        dowodowy. Rdzeń w \path{letsgolegacy.secondkey}, rozszerzenia fazy 02 i 03 w
        \path{secondkey-capture}, \path{secondkey-sandbox} i \path{secondkey-portal}.
  \item \textbf{Portcullis} --- bramka na pull request: deterministyczne analyzery Roslyn
        uruchamiane na zmienionych liniach, z wynikiem w SARIF i plikiem baseline
        (\path{letsgolegacy.portcullis}).
  \item \textbf{secondkey-standards} --- standardy architektury i migracji spisane raz w
        Markdown i dostarczane z jednego źródła agentowi migracyjnemu (jako skill) oraz
        bramce (jako konfiguracja analyzerów).
  \item \textbf{legacy-risk-scan} --- strona, która z pliku projektu (\repopath{.csproj},
        \repopath{packages.config}, \repopath{composer.json}) liczy w przeglądarce raport
        końca wsparcia (EOL) i znanych podatności.
  \item \textbf{bench nopCommerce} --- publiczny okaz: nopCommerce 3.90 na .NET Framework,
        na którym cały łańcuch jest uruchamiany w CI i na którym będzie oceniona migracja
        wykonana przez agenta.
\end{itemize}

Komponenty wymieniają się wyłącznie plikami: nagraniem \repopath{*.skcap}, kontraktem
\repopath{contract.yaml}, przebiegiem \repopath{*.skrun}, werdyktem
\repopath{verdict.json}, logami bramki \repopath{*.sarif} i pakietem dowodowym. Żaden
komponent nie sięga do wnętrza innego, więc każdy da się wymienić, jeśli czyta i pisze te
same artefakty (rozdział~\ref{sec:architektura}).

\zrodla{\path{letsgolegacy.secondkey/README.md},
\path{letsgolegacy.secondkey/docs/WORKPLAN.md},
\path{letsgolegacy.portcullis/docs/WORKPLAN.md},
\path{letsgolegacy.secondkey-standards/README.md},
\path{letsgolegacy.legacy-risk-scan/README.md},
\path{letsgolegacy.secondkey-nopcommerce-bench/README.md}}

\subsection{Nazewnictwo}\label{sec:wprowadzenie-nazwy}

Źródła nazywają cały łańcuch „Second Key” i numerują jego komponenty od C0 do C12 ---
także te, które mieszkają poza rdzeniem: standardy to komponent C4, a Portcullis to
komponent C5, „bramka Second Key”. Ten dokument przyjmuje nazwę nadrzędną
\textbf{letsgolegacy} dla całego frameworku: metody i wszystkich repozytoriów
\repopath{letsgolegacy.*}, łącznie z legacy-risk-scan i benchem, które nie są
komponentami łańcucha. Nazwa \textbf{Second Key} oznacza tu łańcuch weryfikacji i jego
rdzeń (CLI \repopath{sk}, formaty artefaktów), a \textbf{Portcullis} --- bramkę.

Nazwy w kodzie i w formatach się nie zmieniają: przestrzenie nazw \repopath{SecondKey.*},
polecenie \repopath{sk}, plik konfiguracji \repopath{secondkey.yaml}, wersja formatu
kontraktu \path{apiVersion: secondkey/v1}, nazwy repozytoriów i pakietów. Gdy
dokument cytuje źródło, które mówi „Second Key” w szerszym sensie, zaznacza to w
kontekście.

\zrodla{\path{letsgolegacy.secondkey/README.md},
\path{letsgolegacy.portcullis/docs/WORKPLAN.md},
\path{letsgolegacy.secondkey-standards/README.md}}

\subsection{Zakres, wersja i źródła}\label{sec:wprowadzenie-wersja}

Dokument prezentuje publiczne repozytoria frameworku w stanie z tabeli poniżej (gałąź
\repopath{main}). Każda podsekcja kończy się listą plików, z których pochodzi jej treść.
Źródłem prawdy są te pliki: dokumentacja w Markdown, schematy, skrypty, workflowy i kod.
\finding{Gdy ten dokument i źródło się różnią, rację ma źródło}, a dokument należy
poprawić. Nic mechanicznie nie pilnuje, że zmiana w źródle trafi do tego dokumentu --- to
przyjęty koszt ręcznie redagowanej prezentacji; listy źródeł pozwalają sprawdzić każde
twierdzenie osobno.

\begin{tabularx}{\linewidth}{@{}>{\raggedright\arraybackslash}p{6.2cm} l l X@{}}
\toprule
Repozytorium & Commit & Data & Rola \\
\midrule
\path{letsgolegacy.secondkey} & \texttt{fe374c7} & 2026-09-26 & rdzeń Second Key \\
\path{letsgolegacy.portcullis} & \texttt{6a018a2} & 2026-09-26 & bramka C5 \\
\path{letsgolegacy.secondkey-standards} & \texttt{07520bd} & 2026-09-26 & standardy C4 \\
\path{letsgolegacy.secondkey-nopcommerce-bench} & \texttt{e6d4fe9} & 2026-09-26 & publiczny okaz \\
\path{letsgolegacy.legacy-risk-scan} & \texttt{49f8a4c} & 2026-09-26 & raport EOL i CVE \\
\path{letsgolegacy.secondkey-capture} & \texttt{e56068e} & 2026-09-25 & C1, faza 02 \\
\path{letsgolegacy.secondkey-sandbox} & \texttt{75042c6} & 2026-09-25 & C12, faza 02 \\
\path{letsgolegacy.secondkey-portal} & \texttt{e73546f} & 2026-09-25 & C11, faza 03 \\
\bottomrule
\end{tabularx}

Poza zakresem są sprawy nietechniczne i wszystko, czego nie ma w publicznych
repozytoriach. Dokument nie opisuje też pracy w toku na gałęziach, które nie trafiły do
\repopath{main} na wymienionych commitach.

Wersję samego dokumentu --- commit, z którego zbudowano PDF --- podaje nagłówek każdej
strony. PDF buduje workflow \path{.github/workflows/build-docs-pdf.yml} repozytorium
rdzenia przy każdym pushu do \repopath{main}; źródła dokumentu są w
\path{docs/papers/} i \path{docs/diagrams/} (dodatek~\ref{app:ci}).

\zrodla{\path{letsgolegacy.*} na podanych commitach;
\path{architecture-standards/docs/research/00-RESEARCH-DOCUMENTATION.md}}

\subsection{Jak czytać ten dokument}\label{sec:wprowadzenie-czytanie}

\begin{tabularx}{\linewidth}{@{}p{1.9cm} X@{}}
\toprule
Rozdział & Zawartość \\
\midrule
\ref{sec:architektura} & Architektura: łańcuch, artefakty, komponenty C0--C12, zasady projektowe, fazy \\
\ref{sec:procedura} & \textbf{Procedura pracy krok po kroku} \\
\ref{sec:rdzen} & Rdzeń Second Key: projekty, decyzje, testy, CI \\
\ref{sec:formaty} & Formaty artefaktów, w tym pełny język kontraktu i klasy werdyktu \\
\ref{sec:cli} & CLI \texttt{sk}: polecenia, opcje, \texttt{secondkey.yaml}, kody wyjścia \\
\ref{sec:portcullis} & Portcullis: reguły, diff gate, baseline, SARIF, dystrybucja \\
\ref{sec:standardy} & Standardy migracji: 20 reguł, generator, skill, drift check \\
\ref{sec:risk-scan} & legacy-risk-scan \\
\ref{sec:faza02} & Repozytoria fazy 02 i 03: capture, sandbox, portal \\
\ref{sec:bench} & Bench nopCommerce: topologia, skrypty, nagranie, kontrakt, wyniki \\
\ref{sec:bezpieczenstwo} & Bezpieczeństwo i higiena repozytoriów \\
\ref{sec:status} & Status i plan \\
\ref{sec:slownik} & Słownik pojęć \\
\ref{app:klauzule}, \ref{app:ci} & Dodatki: klauzule kontraktu benchu; workflowy CI \\
\bottomrule
\end{tabularx}

Kto chce zacząć pracę, czyta rozdziały~\ref{sec:architektura} i~\ref{sec:procedura};
rozdziały~\ref{sec:rdzen}--\ref{sec:bench} są referencją, do której procedura odsyła.

Konwencje: nazwy techniczne zostają po angielsku, gdy są nazwami własnymi,
identyfikatorami albo ustalonymi terminami (framework, pipeline, pull request, replay,
snapshot); identyfikatory, polecenia i pliki są zapisane krojem kodu, np.
\repopath{sk compare}; wartości z kodu i logów są przepisane bez zmian.

\zrodla{ten dokument}
